% ------------------------------------------------------------------
% P3 appendix -- Proposition \ref{prop:bundle}: the bundled disclosure
% dial has an interior constrained optimum; decoupled dials optimize
% at a corner. Full setup, assumptions (A0)-(A5), lemmas, proofs,
% remarks, falsification record, observation, and conjectures.
%
% Include file: no \documentclass, no \section at top level.
% Requires from the main preamble: amsmath, amssymb, amsthm, and the
% amsthm environments  proposition, lemma, assumption, conjecture,
% observation, remark  (plus  proof  from amsthm).
% ------------------------------------------------------------------

\subsection*{Setup}

This appendix formalizes v9~S11 (\texttt{2026-06-09-ifwishlist-toy-model-v9.md}, \S3).
Ground truth for the model: \texttt{runs/numcheck/v9\_dynamic\_two\_channel.py} (functions
\texttt{solve\_unified}, \texttt{type\_perceived}, \texttt{G\_and\_s},
\texttt{G\_and\_s\_decoupled}); independent re-implementation \texttt{v9\_verify\_s11.py}.
Notation follows the anchor paper \S3--\S6 (\texttt{2026-06-10-ifwishlist-anchor-paper-v3.md}):
disclosure dial $k$, threshold share $\theta$ with $K=\lceil(1-\theta)N\rceil$, excludability
$\xi$, bad-coalition prior $\pi_b$.

\paragraph{The activation game $\Gamma_T(k_{\mathrm{type}},k_{\mathrm{piv}})$.}
The model below is a direct formalization of \texttt{solve\_unified} $+$
\texttt{type\_perceived} in \texttt{v9\_dynamic\_two\_channel.py} --- exact up to the idealized
stage-equilibrium set $S$ (the script approximates $S$ on a grid with bisection; see R1.4) ---
written with two separate dials; the \textbf{bundled} game of S11 is the diagonal
$k_{\mathrm{type}}=k_{\mathrm{piv}}=k$, the \textbf{decoupled} game is the full square
$(k_{\mathrm{type}},k_{\mathrm{piv}})\in[0,1]^2$ (the script's \texttt{G\_and\_s\_decoupled}).

\paragraph{M1 (population, coalition, clearing).}
$N$ agents are attached to one coalition of unknown type $\psi\in\{G,B\}$ with prior
$\Pr(\psi=B)=\pi_b$. Time runs in discrete periods indexed by the remaining horizon
$h=T,T-1,\dots,1$; the state is $(m,h)$ where $m$ $=$ number of agents committed so far. The
coalition \textbf{clears} the first time $m\ge K$.

\paragraph{M2 (payoffs and their perceived, type-blended versions).}
True terminal payoffs: a committer pays $\kappa>0$; if the coalition clears she receives
$v_\psi$ ($v_G=V>0$, $v_B=-L<0$); if she commits and the coalition never clears she receives
$-L_p<0$ (exposure loss, not type-blended). A never-committer receives $\xi v_\psi$ if the
coalition clears ($\xi\in(0,1]$ $=$ excludability leak) and $0$ otherwise. Each period an
uncommitted agent is active she pays waiting friction $\rho\ge 0$. The type dial
$k_{\mathrm{type}}$ enters as a reduced-form \textbf{payoff blend} (the script's
\texttt{type\_perceived}): members of a type-$\psi$ coalition act on the \emph{perceived}
payoffs
\begin{gather*}
P_c(k_{\mathrm{type}};\psi)=(1-k_{\mathrm{type}})\,\bar P_c+k_{\mathrm{type}}\,(v_\psi-\kappa),\\
P_n(k_{\mathrm{type}};\psi)=(1-k_{\mathrm{type}})\,\bar P_n+k_{\mathrm{type}}\,\xi v_\psi,
\end{gather*}
where the pooled values are $\bar P_c=\mathbb E-\kappa$, $\bar P_n=\xi\,\mathbb E$, with
$\mathbb E:=(1-\pi_b)V-\pi_b L$ (prior-mean clearing value). $P_c$ is the perceived payoff to
having committed when the coalition clears; $P_n$ the perceived payoff to not having committed
when it clears.

\paragraph{M3 (stage decision problem).}
At state $(m,h)$ with $m<K$ there are $n=N-m$ uncommitted agents. Given a candidate symmetric
commit propensity $q\in[0,1]$, each \emph{other} uncommitted agent commits this period with
probability $a=p_{\mathrm{opp}}q$ ($p_{\mathrm{opp}}\in(0,1)$ $=$ availability). Let
$Y\sim\mathrm{Binom}(n-1,a)$ be the others' commitments. The deciding agent compares
(continuation objects $\Pi,U$ defined in M6; conventions $\Pi(m',h)=1$ for $m'\ge K$,
$\Pi(m',0)=0$ and $U(m',0)=0$ for $m'<K$):
\[
v^w(q)=-\rho+\mathbb E\Big[\mathbf 1\{m+Y\ge K\}\,P_n+\mathbf 1\{m+Y<K\}\,U(m+Y,h-1)\Big],
\]
\begin{multline*}
v^c(q)=-\rho+\mathbb E\Big[\mathbf 1\{m{+}1{+}Y\ge K\}\,P_c\\
+\mathbf 1\{m{+}1{+}Y<K\}\big(P_c\,\Pi(m{+}1{+}Y,h-1)\\
-L_p\,(1-\Pi(m{+}1{+}Y,h-1))\big)\Big],
\end{multline*}
and the commit--wait differential is $D(q):=v^c(q)-v^w(q)$ (a polynomial in $q$, hence
continuous). The \textbf{stage-equilibrium set} is
\[
S=\{q\in(0,1):D(q)=0\}\cup\{0\ \text{if}\ D(0)\le 0\}\cup\{1\ \text{if}\ D(1)\ge 0\};
\]
$S\neq\emptyset$ by continuity of $D$.

\paragraph{M4 (momentum selection --- a model primitive, not a derived object).}
The played propensity is the \emph{largest} stage equilibrium,
$a_{\mathrm{mom}}(m,h)=\max S$ (the v6/v7 ``optimistic/momentum'' selection; the script's
\texttt{max(eqs)}). Its $\lambda$-free justification is the separate S10 strand; here it is
taken as given.

\paragraph{M5 (pivotality leakage / free-ride damping).}
Define the pivotality weight
\[
\mathrm{piv}(m)=\mathrm{clip}\Big(1-\frac{K-m-1}{\max(1,\;p_{\mathrm{opp}}(N-m))},\,0,\,1\Big)
\]
and the free-ride gate $\chi=\mathbf 1\{\xi V>V-\kappa\}$ (free-riding profitable at true good
payoffs; the script computes the gate once from true values and applies it to both types). A
fraction $f=k_{\mathrm{piv}}\,(1-\mathrm{piv}(m))\,\chi$ of would-be movers free-ride, so the
effective propensity is $a_{\mathrm{eff}}(m,h)=a_{\mathrm{mom}}(m,h)\,(1-f)$.

\paragraph{M6 (transition, values, outputs).}
Given $a_{\mathrm{eff}}$, the number of new commitments is
$J\sim\mathrm{Binom}(n,\,p_{\mathrm{opp}}a_{\mathrm{eff}})$, and
\begin{gather*}
\Pi(m,h)=\mathbb E[\mathbf 1\{m+J\ge K\}+\mathbf 1\{m+J<K\}\,\Pi(m+J,h-1)],\\
U(m,h)=a_{\mathrm{eff}}\,v^c(a_{\mathrm{eff}})+(1-a_{\mathrm{eff}})\,v^w(a_{\mathrm{eff}}).
\end{gather*}
The game is solved once with $\psi=G$ perceived payoffs and once with $\psi=B$; the outputs are
\[
G(k_{\mathrm{type}},k_{\mathrm{piv}})=\Pi_G(0,T),\qquad
s(k_{\mathrm{type}},k_{\mathrm{piv}})=1-\Pi_B(0,T),
\]
\[
W(k_{\mathrm{type}},k_{\mathrm{piv}})=(1-\pi_b)\,G\,V-\pi_b\,(1-s)\,L .
\]
Bundled-dial objects are written $G(k):=G(k,k)$ etc.

\subsection*{Assumptions}

\begin{assumption}[A0 --- regularity]\label{ass:A0}
$N\ge 2$; $K\in\{2,\dots,N\}$; $p_{\mathrm{opp}}\in(0,1)$; $0<\kappa<V$; $L>0$; $L_p>0$;
$\rho\ge 0$; $\pi_b\in(0,1)$; $\xi\in(0,1]$.
\end{assumption}

\begin{assumption}[A1 --- pessimistic pooling]\label{ass:A1}
$0\le \mathbb E<\kappa$, where $\mathbb E=(1-\pi_b)V-\pi_b L$. (Pooled with the bad prior,
clearing is not worth the commitment cost; the pooled free-ride payoff $\xi\mathbb E$ is still
nonnegative.)
\end{assumption}

\begin{assumption}[A2 --- free-ride regime]\label{ass:A2}
$\varepsilon:=\xi V-(V-\kappa)>0$, i.e.\ $\xi>1-\kappa/V$, so the gate $\chi=1$. (The anchor's
``channel has bite only when $\xi$ exceeds $\xi^*=1-\kappa/V$''.)
\end{assumption}

\begin{assumption}[A3 --- thin root field]\label{ass:A3}
$K-1\ge p_{\mathrm{opp}}N$: at the empty state the expected number of available movers does not
cover the gap, so $\mathrm{piv}(0)=0$ --- no early mover is perceived pivotal, and disclosure at
the root reveals pure inframarginality.
\end{assumption}

\begin{assumption}[A4 --- pivotal ignition at full type-news]\label{ass:A4}
With $X\sim\mathrm{Binom}(N-1,p_{\mathrm{opp}})$ and $\beta_1=\Pr(X=K-1)$,
$\beta_+=\Pr(X\ge K)$, $\beta_-=\Pr(X\le K-2)$:
\[
(V-\kappa)\,\beta_1\;>\;\varepsilon\,\beta_+\;+\;L_p\,\beta_-\,.
\]
(At true good payoffs, the pivotal-event gain outweighs the free-ride temptation and the
exposure risk. Since $\beta_1>0$ under A0, A4 is satisfiable whenever $\varepsilon$ and $L_p$
are small enough.)
\end{assumption}

\begin{assumption}[A5 --- small waiting friction]\label{ass:A5}
$\rho\,(T-1)<\min\{L_p,\;\kappa-\mathbb E\}$. (Vacuous at $T=1$.)
\end{assumption}

\paragraph{Non-vacuity.}
The assumption set is non-empty and contains an anchor-shaped instance:
$(N,K,p_{\mathrm{opp}},V,\kappa,\xi,\pi_b,L,L_p)=(10,7,0.6,3,1,0.667,0.5,3,0.5)$, any
$\rho\in[0,1)$ at $T=1$. Check: $\mathbb E=0\in[0,1)$ (A1); $\varepsilon=0.001>0$ (A2);
$K-1=6\ge 6=p_{\mathrm{opp}}N$ (A3); $\beta_1=0.2508$, $\beta_+=0.2318$, $\beta_-=0.5174$, so
A4 reads $0.5016>0.001\cdot 0.2318+0.5\cdot 0.5174=0.2589$ \checkmark.

\paragraph{Stated restrictions relative to the v9 numerics}
(full ledger in the Remarks below): the interior-existence and decoupled-corner parts are
proved at horizon $T=1$ (one activation round), where the continuation objects
$\Pi(\cdot,h-1),U(\cdot,h-1)$ vanish identically; the endpoint results $G(0)=G(1)=0$ are proved
for \textbf{every} $T\ge 1$. The stage-equilibrium set is the idealized $S$ above (the script
approximates it on a grid with bisection).

\subsection*{The proposition}

\begin{restatedproposition}[Proposition~\ref{prop:bundle}, restated: bundled-dial interior optimum; decoupled-dial corner]
\leavevmode
\begin{itemize}
\item[(i)] \textbf{Bundled dial: both corners dead, a live interior, every optimum interior.}
\emph{Endpoints (any horizon).} Under A0, A1, A2, A3, A5, for every $T\ge 1$, the bundled game
satisfies
\[
G(0)=0,\ \ s(0)=1,\qquad G(1)=0,\ \ s(1)=1,\quad\text{hence}\quad W(0)=W(1)=0 .
\]
($k=0$: pooled news cannot ignite momentum --- the unique stage equilibrium is $q=0$ at every
state; $k=1$: full pivotality leakage makes the empty state absorbing --- every root-state mover
free-rides.)
\emph{Interior (one round).} If in addition A4 holds and $T=1$, there is a threshold
$\bar k\in(0,1)$ --- the unique root of the affine root-state differential $D_1$ of
Lemma~\ref{lem:ignition} --- such that for every $k\in[\bar k,1)$:
\begin{gather*}
G(k)=\Pr\!\big(\mathrm{Binom}(N,\;p_{\mathrm{opp}}(1-k))\ge K\big)>0,\qquad s(k)=1,\\
W(k)=(1-\pi_b)\,V\,G(k)>0 .
\end{gather*}
Moreover $W:[0,1]\to\mathbb R$ is upper semicontinuous, attains its maximum, and \textbf{every
maximizer lies in the open interval $(0,1)$}.
\item[(ii)] \textbf{Decoupled dials: the optimum is the corner.} Under A0--A5 and $T=1$, the
decoupled welfare satisfies, for all $(k_{\mathrm{type}},k_{\mathrm{piv}})\in[0,1]^2$,
\[
W(k_{\mathrm{type}},k_{\mathrm{piv}})\;\le\;W(1,0)\;=\;(1-\pi_b)\,V\,
\Pr\!\big(\mathrm{Binom}(N,p_{\mathrm{opp}})\ge K\big)\;>\;0,
\]
and the argmax set is \textbf{exactly} $[\bar k,1]\times\{0\}$ (with $\bar k$ from part (i)):
every optimum has \emph{zero} pivotality leakage and \emph{full ignition}, and
$(k_{\mathrm{type}},k_{\mathrm{piv}})=(1,0)$ attains the optimum. In particular every point
with $k_{\mathrm{piv}}>0$ is strictly suboptimal.
\item[(iii)] \textbf{The interior is a constrained optimum created by bundling.} Under A0--A5
and $T=1$: $\sup_{k\in[0,1]}W(k,k)\le W(1,0)$, and the bundled diagonal never attains $W(1,0)$
(every diagonal point has $a_{\mathrm{eff}}<1$). The bundled optimum is interior (part (i))
while the unconstrained two-dial optimum is the corner (part (ii)); the interior $k^*$ is the
constrained optimum of one dial forced to move both channels together --- not a
type-versus-pivotality balance.
\end{itemize}
\end{restatedproposition}

\subsection*{Proof}

Throughout, fix a state $(m,h)$, write $n=N-m$, $g=K-m$ (the gap), and let
$Y\sim\mathrm{Binom}(n-1,p_{\mathrm{opp}}q)$ denote the others' current-period commitments in
the stage problem. Note
$\Pr(Y=0)=(1-p_{\mathrm{opp}}q)^{\,n-1}\ge(1-p_{\mathrm{opp}})^{\,n-1}>0$
for all $q\in[0,1]$, by $p_{\mathrm{opp}}<1$ (A0). $(\star)$

\begin{lemma}[strict monotonicity of the binomial tail]\label{lem:tail}
For integers $n\ge 1$, $1\le r\le n$, the map $F(b)=\Pr(\mathrm{Binom}(n,b)\ge r)$ is
continuous on $[0,1]$, nondecreasing, strictly increasing on $(0,1)$, and $F(b)>0$ for $b>0$;
indeed $F'(b)=n\binom{n-1}{r-1}b^{\,r-1}(1-b)^{\,n-r}$.
\end{lemma}

\begin{proof}
$F(b)=\sum_{j\ge r}\binom nj b^j(1-b)^{n-j}$ is a polynomial. Term-by-term,
\[
\frac{d}{db}\binom nj b^j(1-b)^{n-j}
=\underbrace{n\binom{n-1}{j-1}b^{j-1}(1-b)^{n-j}}_{=:t_j}
-\underbrace{n\binom{n-1}{j}b^{j}(1-b)^{n-j-1}}_{=:u_j},
\]
using $j\binom nj=n\binom{n-1}{j-1}$ and $(n-j)\binom nj=n\binom{n-1}{j}$. Since $u_j=t_{j+1}$
and $u_n=0$, the sum over $j=r,\dots,n$ telescopes to
$t_r=n\binom{n-1}{r-1}b^{r-1}(1-b)^{n-r}$, which is $>0$ on $(0,1)$. Positivity for $b>0$:
$F(b)\ge\Pr(\mathrm{Binom}(n,b)=n)=b^n>0$ when $r\le n$.
\end{proof}

\begin{lemma}[no ignition under pooled news; any horizon, any $k_{\mathrm{piv}}$, both types]\label{lem:noignition}
Assume A0, A5, the bound $\mathbb E<\kappa$ from A1 (the sign restriction $\mathbb E\ge 0$ in
A1 is not used here), and $k_{\mathrm{type}}=0$ (with $k_{\mathrm{piv}}\in[0,1]$ arbitrary,
type $\psi\in\{G,B\}$ arbitrary). Then for every state $(m,h)$ with $m<K$, $h\ge 1$: the unique
stage equilibrium is $q=0$, hence $a_{\mathrm{mom}}=a_{\mathrm{eff}}=0$; and $\Pi(m,h)=0$,
$U(m,h)=-\rho h$. Consequently $G(0,k_{\mathrm{piv}})=0$ and $s(0,k_{\mathrm{piv}})=1$ for
every $T\ge 1$.
\end{lemma}

\begin{proof}
At $k_{\mathrm{type}}=0$ the perceived payoffs are the pooled values for \textbf{both} types:
$P_c=\bar P_c=\mathbb E-\kappa$, $P_n=\bar P_n=\xi\mathbb E$. By $\mathbb E<\kappa$:
$\bar P_c<0$, and $\bar P_c-\bar P_n=(1-\xi)\mathbb E-\kappa<0$ (if $\mathbb E\ge 0$ this is
$\le\mathbb E-\kappa<0$ using $\xi\le 1$; if $\mathbb E<0$ then $(1-\xi)\mathbb E\le 0$ and the
expression is $\le-\kappa<0$). $(\dagger)$

Induct on $h$. \emph{Base $h=0$:} $\Pi(m,0)=0$, $U(m,0)=0=-\rho\cdot 0$ for $m<K$ by
convention. \emph{Step:} assume $\Pi(m',h-1)=0$ and $U(m',h-1)=-\rho(h-1)$ for all $m'<K$.
Substituting the hypothesis into M3 (failure branch of $v^c$ becomes $-L_p$; wait branch on
$\{m+Y<K\}$ becomes $-\rho(h-1)$; the two leading $-\rho$ cancel):
\begin{multline*}
D(q)=\Big[\bar P_c\Pr(Y\ge g-1)-L_p\Pr(Y\le g-2)\Big]\\
-\Big[\bar P_n\Pr(Y\ge g)-\rho(h-1)\Pr(Y\le g-1)\Big].
\end{multline*}
Splitting $\Pr(Y\ge g-1)=\Pr(Y=g-1)+\Pr(Y\ge g)$ and
$\Pr(Y\le g-1)=\Pr(Y=g-1)+\Pr(Y\le g-2)$:
\begin{equation}
\begin{split}
D(q)={}&\big(\bar P_c+\rho(h-1)\big)\Pr(Y=g-1)
+\big(\bar P_c-\bar P_n\big)\Pr(Y\ge g)\\
&-\big(L_p-\rho(h-1)\big)\Pr(Y\le g-2).
\end{split}\tag{L2.1}
\end{equation}
By A5, $\rho(h-1)\le\rho(T-1)<\min\{L_p,\kappa-\mathbb E\}$, so all three coefficients are
negative: $\bar P_c+\rho(h-1)=\mathbb E-\kappa+\rho(h-1)<0$; $\bar P_c-\bar P_n<0$ by
$(\dagger)$; $-(L_p-\rho(h-1))<0$. Hence $D(q)\le 0$ for all $q$. Strictness: if $m\le K-2$
then $g-2\ge 0$ so $\Pr(Y\le g-2)\ge\Pr(Y=0)>0$ by $(\star)$ and the third term is strictly
negative; if $m=K-1$ then $g-1=0$ so $\Pr(Y=g-1)\ge\Pr(Y=0)>0$ by $(\star)$ and the first term
is strictly negative. So $D(q)<0$ for \textbf{all} $q\in[0,1]$: $q=0\in S$ (since $D(0)<0$),
no $q\in(0,1)$ is in $S$ (needs $D=0$), and $1\notin S$ (needs $D(1)\ge 0$). Thus
$a_{\mathrm{mom}}=\max S=0$, so $a_{\mathrm{eff}}=0$ regardless of $k_{\mathrm{piv}}$ (the
damping multiplies zero). The transition is then degenerate at $J=0$:
$\Pi(m,h)=\Pi(m,h-1)=0$, and $U(m,h)=v^w(0)=-\rho+U(m,h-1)=-\rho h$ (at $h=1$ the
wait-no-clear branch is $0$, consistent with the convention). The induction closes. Finally
$G(0,k_{\mathrm{piv}})=\Pi_G(0,T)=0$ and $\Pi_B(0,T)=0$ gives $s=1$.
\end{proof}

\begin{lemma}[full pivotality leakage freezes the root; any horizon, any $k_{\mathrm{type}}$, both types]\label{lem:freeze}
Assume A0, A2, A3 and $k_{\mathrm{piv}}=1$ (with $k_{\mathrm{type}}\in[0,1]$, $\psi$
arbitrary). Then $a_{\mathrm{eff}}(0,h)=0$ for every $h$, the state $m=0$ is absorbing, and
$\Pi(0,T)=0$ for every $T\ge 1$. Consequently $G(k_{\mathrm{type}},1)=0$ and
$s(k_{\mathrm{type}},1)=1$.
\end{lemma}

\begin{proof}
At $m=0$: the gap is $K$ and $n=N$, so
$\mathrm{piv}(0)=\mathrm{clip}\big(1-\frac{K-1}{\max(1,\,p_{\mathrm{opp}}N)}\big)$. If
$p_{\mathrm{opp}}N\ge 1$ the ratio is $\ge 1$ by A3 ($K-1\ge p_{\mathrm{opp}}N$); if
$p_{\mathrm{opp}}N<1$ the denominator is $1$ and the ratio is $K-1\ge 1$ by $K\ge 2$ (A0).
Either way $\mathrm{piv}(0)=0$. By A2 the gate $\chi=1$. Hence $f=1\cdot(1-0)\cdot 1=1$ and
$a_{\mathrm{eff}}(0,h)=a_{\mathrm{mom}}(0,h)\cdot 0=0$ --- whatever $a_{\mathrm{mom}}$,
$k_{\mathrm{type}}$, and $\psi$ are. The transition from $(0,h)$ is degenerate at $J=0$, so by
induction on $h$, $\Pi(0,T)=\Pi(0,0)=0$ (using $K\ge 2>0$). Both types are frozen at the root,
so $G(\cdot,1)=0$ and $1-s(\cdot,1)=\Pi_B(0,T)=0$.
\end{proof}

\begin{lemma}[interior ignition at $T=1$]\label{lem:ignition}
Assume A0--A5 and $T=1$, type $G$, bundled dial $k_{\mathrm{type}}=k_{\mathrm{piv}}=k$. Define
\[
D_1(k):=P_c(k;G)\,\beta_1-\big(P_n(k;G)-P_c(k;G)\big)\,\beta_+-L_p\,\beta_-\,,
\]
with $\beta_1,\beta_+,\beta_-$ as in A4. Then: (a) $D_1$ is the commit--wait differential at
the root state $(0,1)$ evaluated at $q=1$; (b) $D_1$ is affine and strictly increasing in $k$,
with $D_1(0)<0<D_1(1)$, so there is a unique $\bar k\in(0,1)$ with $D_1(k)\ge 0$ exactly on
$[\bar k,1]$; (c) for every $k\in[\bar k,1]$, $a_{\mathrm{mom}}(0,1)=1$ and
$a_{\mathrm{eff}}(0,1)=1-k$; (d) for every $k\in[\bar k,1)$,
$G(k)=\Pr(\mathrm{Binom}(N,p_{\mathrm{opp}}(1-k))\ge K)>0$.
\end{lemma}

\begin{proof}
(a) At $T=1$ the continuations vanish ($\Pi(\cdot,0)=0$, $U(\cdot,0)=0$), so at
$(m,h)=(0,1)$, $g=K$, the formula-(L2.1)-style computation gives, for general $q$,
\begin{equation}
D(q;k)=P_c(k)\Pr(Y=K-1)-\big(P_n(k)-P_c(k)\big)\Pr(Y\ge K)-L_p\Pr(Y\le K-2),\tag{L4.1}
\end{equation}
where now $P_c(k)=P_c(k;G)$, $P_n(k)=P_n(k;G)$ and $Y\sim\mathrm{Binom}(N-1,p_{\mathrm{opp}}q)$.
At $q=1$ the three probabilities are $\beta_1,\beta_+,\beta_-$, giving $D(1;k)=D_1(k)$.

(b) $P_c(\cdot)$ and $P_n(\cdot)$ are affine in $k$ (M2), so $D_1$ is affine. At $k=1$:
$P_c=V-\kappa$, $P_n-P_c=\xi V-(V-\kappa)=\varepsilon$, so
$D_1(1)=(V-\kappa)\beta_1-\varepsilon\beta_+-L_p\beta_->0$ by A4. At $k=0$ the payoffs are the
pooled ones, and each term of (L4.1) is $\le 0$ by $(\dagger)$ of
Lemma~\ref{lem:noignition} ($\bar P_c<0$, $\bar P_c-\bar P_n<0$), with the third strictly
negative because $\beta_-\ge(1-p_{\mathrm{opp}})^{N-1}>0$ ($(\star)$, using $K\ge 2$); so
$D_1(0)<0$. An affine function with $D_1(0)<0<D_1(1)$ has strictly positive slope, hence a
unique root $\bar k\in(0,1)$ and $\{k:D_1(k)\ge 0\}=[\bar k,1]$.

(c) For $k\in[\bar k,1]$, $D(1;k)=D_1(k)\ge 0$, so $q=1\in S$ and
$a_{\mathrm{mom}}(0,1)=\max S=1$. By A3 and A2 (as in Lemma~\ref{lem:freeze})
$\mathrm{piv}(0)=0$, $\chi=1$, so $f=k$ and $a_{\mathrm{eff}}=1-k$.

(d) For $k\in[\bar k,1)$, $a_{\mathrm{eff}}=1-k\in(0,1]$, and at $T=1$ the clearing
probability is exactly one binomial tail over all $n=N$ uncommitted agents (M6 with
$\Pi(\cdot,0)=0$):
$G(k)=\Pi_G(0,1)=\Pr\big(\mathrm{Binom}(N,\,p_{\mathrm{opp}}(1-k))\ge K\big)$, which is $>0$
by Lemma~\ref{lem:tail} since $p_{\mathrm{opp}}(1-k)>0$ and $K\le N$.
\end{proof}

\begin{lemma}[bad coalitions never ignite at $T=1$, at any dial setting]\label{lem:bad}
Assume A0, A1 and $T=1$, type $B$, any $(k_{\mathrm{type}},k_{\mathrm{piv}})\in[0,1]^2$. (A2
is not needed: the stage equilibrium below is already $q=0$, so the conclusion holds whatever
the value of the gate $\chi$.) Then the unique stage equilibrium at $(0,1)$ is $q=0$, so
$\Pi_B(0,1)=0$ and $s(k_{\mathrm{type}},k_{\mathrm{piv}})=1$.
\end{lemma}

\begin{proof}
Write $t=k_{\mathrm{type}}$. The bad-type perceived payoffs are
$P_c(t;B)=(1-t)\bar P_c+t(-L-\kappa)$ and $P_n(t;B)=(1-t)\bar P_n+t(-\xi L)$. Both blend
components of $P_c(\cdot;B)$ are negative ($\bar P_c=\mathbb E-\kappa<0$ by A1;
$-L-\kappa<0$), so $P_c(t;B)<0$ for all $t$. Also
\[
P_c(t;B)-P_n(t;B)=(1-t)\big((1-\xi)\mathbb E-\kappa\big)+t\big(-(1-\xi)L-\kappa\big)<0,
\]
since the first bracket is $<0$ by $(\dagger)$ and the second is $<0$ ($\xi\le 1$, $L>0$,
$\kappa>0$). Substituting into (L4.1) (which holds for either type at $T=1$): all three terms
are $\le 0$ and the third is strictly negative ($\Pr(Y\le K-2)\ge\Pr(Y=0)>0$ by $(\star)$ and
$K\ge 2$). Hence $D(q)<0$ for all $q$, $a_{\mathrm{mom}}=0$, $a_{\mathrm{eff}}=0$,
$\Pi_B(0,1)=0$.
\end{proof}

\begin{lemma}[upper semicontinuity and attainment, $T=1$, bundled]\label{lem:usc}
Assume A0--A3, A5 and $T=1$. Then $k\mapsto a_{\mathrm{mom}}(0,1)(k)$ is upper semicontinuous
on $[0,1]$; hence
$W(k)=(1-\pi_b)V\cdot\Pr(\mathrm{Binom}(N,p_{\mathrm{opp}}\,a_{\mathrm{mom}}(k)(1-k))\ge K)$
is upper semicontinuous and attains its maximum on $[0,1]$.
\end{lemma}

\begin{proof}
By Lemma~\ref{lem:bad}, $s\equiv 1$, so $W(k)=(1-\pi_b)V\,G(k)$ with
$G(k)=F\big(p_{\mathrm{opp}}\,a_{\mathrm{eff}}(k)\big)$, $F$ the (continuous, nondecreasing)
binomial tail of Lemma~\ref{lem:tail} and $a_{\mathrm{eff}}(k)=a_{\mathrm{mom}}(k)(1-k)$
(Lemma~\ref{lem:ignition}(c)'s damping form, valid at the root for all $k$ by A2--A3).
$D(q;k)$ in (L4.1) is jointly continuous (polynomial in $q$, affine in $k$). Let $k_j\to k$
and pass to a subsequence with
$q_j:=a_{\mathrm{mom}}(k_j)\to\alpha=\limsup_j a_{\mathrm{mom}}(k_j)$. Each $q_j\in S(k_j)$
falls in one of three cases, one of which holds along a further subsequence: ($q_j=0$,
$D(0;k_j)\le 0$): then $\alpha=0\le a_{\mathrm{mom}}(k)$ since $S(k)\subseteq[0,1]$ is
nonempty. ($q_j=1$, $D(1;k_j)\ge 0$): then $D(1;k)\ge 0$ by continuity, so $1\in S(k)$ and
$a_{\mathrm{mom}}(k)=1\ge\alpha$. ($q_j\in(0,1)$, $D(q_j;k_j)=0$): then $D(\alpha;k)=0$ by
joint continuity; if $\alpha\in(0,1)$ then $\alpha\in S(k)$; if $\alpha\in\{0,1\}$ then
$D(\alpha;k)=0$ also certifies $\alpha\in S(k)$ (the endpoint inequalities hold with
equality). Either way $a_{\mathrm{mom}}(k)\ge\alpha$. So $a_{\mathrm{mom}}$ is u.s.c.; then
$a_{\mathrm{eff}}$ is u.s.c.\ (product of a nonnegative u.s.c.\ function and a continuous
nonnegative function), and $W=(1-\pi_b)V\cdot F(p_{\mathrm{opp}}\cdot)$ composed with
$a_{\mathrm{eff}}$ is u.s.c.\ ($F$ continuous and nondecreasing). An u.s.c.\ function on a
compact set attains its maximum.
\end{proof}

\begin{proof}[Proof of Proposition~\ref{prop:bundle}(i)]
\emph{Endpoints.} $G(0)=0$ and $s(0)=1$: Lemma~\ref{lem:noignition} with $k_{\mathrm{piv}}=0$
(any $T$). $G(1)=0$ and $s(1)=1$: Lemma~\ref{lem:freeze} with $k_{\mathrm{type}}=1$ (any $T$).
Then by M6, $W(0)=W(1)=(1-\pi_b)\cdot 0\cdot V-\pi_b\cdot 0\cdot L=0$.

\emph{Interior, $T=1$.} Lemma~\ref{lem:ignition} gives $\bar k\in(0,1)$ and, for
$k\in[\bar k,1)$, $G(k)=\Pr(\mathrm{Binom}(N,p_{\mathrm{opp}}(1-k))\ge K)>0$;
Lemma~\ref{lem:bad} gives $s(k)=1$, so $W(k)=(1-\pi_b)V\,G(k)>0$ (using $\pi_b<1$, $V>0$).

\emph{Every maximizer interior.} By Lemma~\ref{lem:usc} the maximum is attained; its value is
$\ge W(\bar k)>0=W(0)=W(1)$, so neither endpoint is a maximizer; every maximizer lies in
$(0,1)$.
\end{proof}

\begin{proof}[Proof of Proposition~\ref{prop:bundle}(ii)]
At $T=1$, by Lemma~\ref{lem:bad}, $s\equiv 1$ on $[0,1]^2$, so
$W(k_{\mathrm{type}},k_{\mathrm{piv}})=(1-\pi_b)V\,F\big(p_{\mathrm{opp}}\,
a_{\mathrm{eff}}(k_{\mathrm{type}},k_{\mathrm{piv}})\big)$ with
$F(b)=\Pr(\mathrm{Binom}(N,b)\ge K)$ and, at the root (A2, A3, as in
Lemma~\ref{lem:freeze}),
\[
a_{\mathrm{eff}}(k_{\mathrm{type}},k_{\mathrm{piv}})
=a_{\mathrm{mom}}(k_{\mathrm{type}})\,(1-k_{\mathrm{piv}}),
\]
where $a_{\mathrm{mom}}(k_{\mathrm{type}})$ depends only on $k_{\mathrm{type}}$ (the stage
problem M3 involves the perceived payoffs only; the damping M5 is applied after).

\emph{Upper bound and attainment.} $a_{\mathrm{eff}}\le 1$ always, and $F$ is nondecreasing
(Lemma~\ref{lem:tail}), so $W\le(1-\pi_b)V\,F(p_{\mathrm{opp}})$. By
Lemma~\ref{lem:ignition}(b)--(c) (whose ignition part uses only the type dial),
$a_{\mathrm{mom}}(k_{\mathrm{type}})=1$ exactly for $k_{\mathrm{type}}\in[\bar k,1]$; in
particular $a_{\mathrm{mom}}(1)=1$ (A4). Hence $a_{\mathrm{eff}}(1,0)=1$ and
$W(1,0)=(1-\pi_b)V\,F(p_{\mathrm{opp}})$, which is $>0$ by Lemma~\ref{lem:tail}
($p_{\mathrm{opp}}>0$, $K\le N$). So the upper bound is attained at $(1,0)$.

\emph{Exact argmax set.} A point $(k_{\mathrm{type}},k_{\mathrm{piv}})$ attains the bound iff
$F(p_{\mathrm{opp}}a_{\mathrm{eff}})=F(p_{\mathrm{opp}})$. Since $F$ is strictly increasing on
$(0,1)$ (Lemma~\ref{lem:tail}) and $0<p_{\mathrm{opp}}<1$, this holds iff
$a_{\mathrm{eff}}=1$, i.e.\ iff $a_{\mathrm{mom}}(k_{\mathrm{type}})=1$ \textbf{and}
$k_{\mathrm{piv}}=0$, i.e.\ iff
$(k_{\mathrm{type}},k_{\mathrm{piv}})\in[\bar k,1]\times\{0\}$. (If $a_{\mathrm{eff}}=0$ the
value is $F(0)=0<W(1,0)$; if $a_{\mathrm{eff}}\in(0,1)$ the value is strictly below the
bound.) In particular any $k_{\mathrm{piv}}>0$ forces
$a_{\mathrm{eff}}\le 1-k_{\mathrm{piv}}<1$, hence strict suboptimality.
\end{proof}

\begin{proof}[Proof of Proposition~\ref{prop:bundle}(iii)]
The diagonal is a subset of the square, so $\sup_k W(k,k)\le\max W=W(1,0)$. No diagonal point
attains it: attainment requires $a_{\mathrm{eff}}=1$, i.e.\ $a_{\mathrm{mom}}(k)=1$ and $k=0$
simultaneously; but $a_{\mathrm{mom}}(0)=0$ by Lemma~\ref{lem:noignition} (so $k=0$ gives
$a_{\mathrm{eff}}=0$), and $k>0$ gives $a_{\mathrm{eff}}\le 1-k<1$. Combined with part (i)
(every bundled maximizer interior) and part (ii) (every decoupled maximizer at
$k_{\mathrm{piv}}=0$, $k_{\mathrm{type}}\in[\bar k,1]$), the bundled interior optimum is
strictly below the decoupled corner optimum: it is the constrained optimum of the yoked dial,
not a channel balance.
\end{proof}

\subsection*{Remarks}

\begin{remark}[R1 --- what was restricted relative to the v9 numerics]\label{rem:ledger}
\leavevmode
\begin{enumerate}
\item \emph{Horizon.} The interior-existence step (Lemma~\ref{lem:ignition}) and the
decoupled-corner result (part (ii)) are proved at $T=1$, where the continuation objects
vanish. The endpoint results $G(0)=G(1)=0$ (Lemmas~\ref{lem:noignition}--\ref{lem:freeze})
hold for \textbf{every} $T\ge 1$, including the script's $T=10$. The script's run is the
$T=10$ game; see the conjecture block below (C1 open; C2 falsified).
\item \emph{Assumption regime vs.\ the anchor calibration.} The anchor calibration
($\pi_b=0.2$, $V=L=3$, $\kappa=1$) has $\mathbb E=1.8>\kappa$, violating A1. It also violates
A4: with the anchor's $L_p=3$, the right side of A4 is
$\varepsilon\beta_++L_p\beta_-=0.001\cdot 0.2318+3\cdot 0.5174=1.5524$, far above the left
side $(V-\kappa)\beta_1=2\cdot 0.2508=0.5016$. And A5 fails automatically once A1 does
($\kappa-\mathbb E=-0.8<0$, so no $\rho\ge 0$ satisfies it at $T=10$). Of A0--A5 the anchor
thus satisfies only A0, A2, A3: no interior result (Lemma~\ref{lem:ignition}, part (i)'s
interior clause, part (ii)) is instantiated at the anchor --- see R3 for what this does to the
mapping. Its $G(0)\approx 0$ is produced by a different sufficient mechanism --- large exposure
risk $L_p=3$ plus a thin field $p_{\mathrm{opp}}=0.6$ (even at $q=1$ the root failure
probability is large, so the $-L_p$ term dominates and the high stage equilibrium fails) ---
which we have \textbf{not} formalized. A1 (pessimistic pooling, $\mathbb E<\kappa$) is the
clean sufficient route to the same endpoint, and is the economically natural one (``masking
pools you with the bad prior, and the pool is not worth the commitment cost''). The conditions
here are sufficient, not necessary, and not the anchor calibration's own mechanism at $k=0$.
The $k=1$ endpoint mechanism (Lemma~\ref{lem:freeze}), by contrast, \textbf{is} exactly the
anchor's: in the $T=10$ run, $G(1)=0$ arises because $f=1$ at $m=0$
($K-1=6\ge 6=p_{\mathrm{opp}}N$, A3 holds with equality at the anchor) --- the same
root-freezing as our proof.
\item \emph{Equilibrium selection.} The optimistic selection $a_{\mathrm{mom}}=\max S$ is a
model primitive (M4, the v6/v7 convention; its $\lambda$-free license is the S10/FMP \S6.4
strand, not re-proved here). Lemma~\ref{lem:ignition}'s interior ignition shows $q=1$
\emph{is} a stage equilibrium and relies on the max selection to conclude it is played. Under
a pessimistic selection ($\min S$) the interior would die at $T=1$ ($D(0)<0$ always, so
$0\in S$ always); the inverted-U is a property of the momentum-selected dynamics, exactly as
v9 frames it.
\item \emph{Idealized equilibrium set.} The script approximates $S$ on a $201$-point grid with
bisection and $10^{-12}$ tolerances; we use the exact $S$. No result depends on the
approximation.
\item \emph{Reduced-form channels.} The type channel is a payoff blend, not Bayesian updating
from a signal, and the free-ride fraction is the script's
$k_{\mathrm{piv}}(1-\mathrm{piv})\chi$ form --- both inherited verbatim from the ground-truth
script (its own header flags them as ``honestly a unified-but-reduced model''). P3 formalizes
that reduced model, not the deeper S12 information-theoretic channel.
\end{enumerate}
\end{remark}

\begin{remark}[R2 --- a fully verified instance]\label{rem:instance}
$(N,K,p_{\mathrm{opp}},V,\kappa,\xi,\pi_b,L,L_p,T)=(10,7,0.6,3,1,0.667,0.5,3,0.5,1)$ satisfies
A0--A5 (arithmetic in the Setup above). Here $\bar k$ solves $0.98402\,k=0.74130$, i.e.\
$\bar k\approx 0.7533$. Cross-check against the ground-truth solver run at these parameters
with $T=1$ (sanity check only; no proof step uses it): $G(k)=0$ for $k\in\{0,0.2,\dots,0.75\}$,
$G(0.76)=1.03\times10^{-4}>0$, $G(0.9)=2.9\times10^{-7}>0$, $G(1)=0$, $s\equiv 1$; decoupled
$11\times 11$ grid argmax value $0.57342=(1-\pi_b)V\,F(p_{\mathrm{opp}})$ exactly as
part (ii) predicts, reported at $(k_{\mathrm{type}},k_{\mathrm{piv}})=(0.8,0.0)$ --- the
smallest grid point of the predicted flat argmax set
$[\bar k,1]\times\{0\}=[0.7533,1]\times\{0\}$ (the script's tie-break keeps the first
maximizer found).
\end{remark}

\begin{remark}[R3 --- map to the numerical anchors (v9 S11, $T=10$, $\theta=0.30$, $\xi=0.667$, $\pi_b=0.2$)]\label{rem:anchormap}
$\theta=0.30\mapsto K=7$ --- A3 holds with equality ($K-1=6=p_{\mathrm{opp}}N$).
$\xi=0.667\mapsto$ A2 with $\varepsilon=0.001$ (the anchor sits just inside the free-ride
regime, $\xi^*=2/3$). ``$G(0)=0$: masking cannot ignite momentum'' $\mapsto$
Lemma~\ref{lem:noignition}'s no-ignition conclusion ($a\equiv 0$ at every state; sufficient
condition A1$+$A5, see R1.2). ``$G(1)=0$: full leakage fully unravels'' $\mapsto$
Lemma~\ref{lem:freeze} (root absorbing; the anchor satisfies the exact hypotheses A2$+$A3).
``Interior $G(0.70)=0.997$, $W=+2.39$'' $\mapsto$ Lemma~\ref{lem:ignition}'s live band
$[\bar k,1)$ --- illustrative only, NOT an instantiation: the anchor violates A1, A4, and A5
(R1.2), so Lemma~\ref{lem:ignition}'s hypotheses do not hold there and the anchor's interior
is not an instance of the proposition. Among the anchor facts, only the $G(1)=0$ endpoint
(Lemma~\ref{lem:freeze}, hypotheses A2$+$A3, which the anchor meets exactly) is a genuine
instance; the $G(0)=0$ endpoint and the interior are qualitative matches produced by
mechanisms we have not formalized. The genuine A0--A5 instance with a live interior is R2's
($L_p=0.5$). At $T=1$ magnitudes are tiny ($G\sim 10^{-4}$) because a single round has no
momentum amplification, while the anchor's $T=10$ chain amplifies the same ignition into
$G\approx 1$. ``Decoupled corner $(k_{\mathrm{type}},k_{\mathrm{piv}})=(1,0)$ at every $\xi$''
$\mapsto$ part (ii): the corner $(1,0)$ attains the optimum for every $\xi$ satisfying A2
(the proof nowhere fixes $\xi$ beyond A2/A4); for $\xi\le\xi^*$ the gate is off, $f\equiv 0$,
$k_{\mathrm{piv}}$ is payoff-irrelevant, and $k_{\mathrm{piv}}=0$ is (weakly) optimal
trivially. ``Interior $W=+2.39$ beats both corners by the whole range'' $\mapsto$
Proposition~\ref{prop:bundle}(i) final clause: $\max W>0=W(0)=W(1)$.
\end{remark}

\subsection*{Conjectures, a falsification record, and an observation (R4)}

The following records the precise open statements and what would close them, together with the
falsification record of the original general-horizon decoupled-corner conjecture.

\begin{restatedconjecture}[Conjecture~\ref{conj:interiorT} (C1), restated: general-horizon interior]
Under A0--A3, A5, and a strengthened ignition condition, for every $T\ge 1$ there exists
$k\in(0,1)$ with $G(k)>0$ in the bundled game.
\end{restatedconjecture}

\emph{Status:} proved at $T=1$ (Lemma~\ref{lem:ignition}); numerically corroborated at $T=10$
at the anchor (live interior with $G(0.70)=0.997$ --- corroboration only, since the anchor lies
outside A0--A5, R1.2). \emph{Obstruction:} for $T\ge 2$ the wait branch of $D(1)$ contains the
endogenous option value $U(m+Y,T-1)$ of staying uncommitted; the crude bound $U\le P_n$ is too
loose (it makes the wait branch dominate for any $\varepsilon>0$). Closing C1 needs a sharper
upper bound on the waiting option value at momentum-selected play --- e.g.,
$U(m',h)\le P_n\Pi(m',h)-c\,(1-\Pi(m',h))$ for a positive constant $c$ --- or a direct
monotone-induction argument that ignition at the root only improves at deeper states (where
the gap is smaller and $\mathrm{piv}>0$ weakens the damping).

\begin{remark}[A falsified conjecture: C2 --- general-horizon decoupled corner]\label{rem:falsifiedC2}
As originally conjectured, C2 read: under A0--A5, for every $T\ge 1$ the decoupled argmax has
$k_{\mathrm{piv}}=0$ and is attained at $(1,0)$. This is false. The counterexample hunter
exhibited, inside A0--A5,
\begin{gather*}
(N,K,p_{\mathrm{opp}},V,\kappa,\xi,\pi_b,L,L_p,\rho)\\
=(8,4,0.29,4.84,2.00,0.958,0.32,9.625,0.408,0.0126)
\end{gather*}
--- all of A0--A5 verified, A4: $0.6160>0.4770$ --- where the ground-truth solver gives
$W(k_{\mathrm{type}}{=}1,k_{\mathrm{piv}}{=}0.2)=0.9903>W(1,0)=0.7629$, identically at
$T=2,5,10$, with $s\equiv 1$ --- a pure good-coalition effect: at the full-disclosure corner,
leakage $k_{\mathrm{piv}}\colon 0\to 0.2$ strictly RAISES the activation gain (the hunter
reports the gain rising from $0.232$ to $0.296$; our re-extraction from
\texttt{v9\_dynamic\_two\_channel.py} gives $W(1,0.2)=0.99025$, $W(1,0)=0.76286$, with
$G\colon 0.2318\to 0.3009$). Confirmed by \texttt{v9\_verify\_s11.W\_general} to 5 decimals.
Per the hunter, the failure is GENERIC: \texttt{6/8 fresh random A0-A5 draws violate C2 at
T=10 (k\_piv in \{0.1-0.3\} strictly beats every k\_piv=0 point)}, and $(1,0)$ is not even the
argmax of the $k_{\mathrm{piv}}=0$ line at $T\ge 2$ ($k_{\mathrm{type}}=0.9$ beats
$k_{\mathrm{type}}=1$ in the instance above) --- we re-verified the line claim:
$W(0.9,0)=0.7974>0.7629=W(1,0)$ at $T=10$, and $W(0.8,0)=0.8284>W(1,0)$ already at $T=2$.

\emph{Replication note (recorded for completeness, not as a hedge --- the falsification stands
on the re-verified numbers above):} one hunter figure did not reproduce, the ``max over the
entire $k_{\mathrm{piv}}=0$ line $=0.9378$''. On an 11-point $k_{\mathrm{type}}$ grid the line
max is $0.8284$ (so $W(1,0.2)$ does beat every coarse-grid $k_{\mathrm{piv}}=0$ point), but
refining the line yields $W(0.85,0)=1.0313$ and $W(0.8625,0)=1.0788$, both above $W(1,0.2)$,
and a $21\times 21$ scan with local refinement places this instance's global argmax at
$(0.8625,0)$ --- at $k_{\mathrm{piv}}=0$ after all. What is decisively falsified is therefore
C2's attainment at $(1,0)$ and the part-(ii) comparative static ``every $k_{\mathrm{piv}}>0$
point is strictly suboptimal''; whether the weaker ``some $k_{\mathrm{piv}}=0$ point is
optimal'' survives is C2$'$ below (Conjecture~\ref{conj:decoupledT}; the hunter's 6/8
line-dominance draws carry the same coarse-grid caveat).

\emph{What survives:} (a) the decoupled corner is a THEOREM at $T=1$
(Proposition~\ref{prop:bundle}(ii), proved above, untouched); (b) at the v9 anchor calibration
at $T=10$ it is a numerical fact --- re-verified against
\texttt{v9\_dynamic\_two\_channel.best\_2d} ($11\times 11$): argmax $(1.00,0.00)$ with
$W^*=+2.4000,\ +2.3995,\ +2.3855$ at $\xi=0.667,\ 0.70,\ 0.75$. (The anchor lies outside
A0--A5, R1.2, so this corroborates rather than instantiates.)

\emph{Mechanism (per the hunter; consistent with our re-runs):} at $T\ge 2$ a leakier future
cuts the waiting option value and re-ignites mid-states --- exactly the continuation channel
named in C1's obstruction as what breaks the $T=1$ factorization
$a_{\mathrm{eff}}=a_{\mathrm{mom}}(k_{\mathrm{type}})(1-k_{\mathrm{piv}})$.
\end{remark}

\begin{restatedobservation}[Observation~\ref{obs:leakagehelps} (O1), restated: positive leakage can strictly help at $T\ge 2$]
There exist instances satisfying A0--A5, horizons $T\ge 2$, and dials
$k_{\mathrm{type}}\in[0,1]$, $k_{\mathrm{piv}}>0$ such that
$W(k_{\mathrm{type}},k_{\mathrm{piv}})>W(k_{\mathrm{type}},0)$ with $s\equiv 1$ at both
points: positive pivotality leakage is strictly welfare-improving, through the good-coalition
continuation/waiting-option channel. Verified instance: the C2 counterexample
(Remark~\ref{rem:falsifiedC2}) at $k_{\mathrm{type}}=1$, $k_{\mathrm{piv}}=0.2$, where $W$
rises $0.7629\to 0.9902$ for every $T\in\{2,5,10\}$. At $T=1$ this is impossible
(Proposition~\ref{prop:bundle}(ii): $a_{\mathrm{eff}}$ is nonincreasing in $k_{\mathrm{piv}}$
and $F$ is nondecreasing), so the phenomenon is strictly dynamic.
\end{restatedobservation}

\begin{restatedconjecture}[Conjecture~\ref{conj:decoupledT} (C2$'$), restated: narrowed residual conjecture]
Under A0--A5, for every $T\ge 1$ the decoupled welfare attains its maximum at some point with
$k_{\mathrm{piv}}=0$ --- with no claim that $(1,0)$ attains it, and no pointwise dominance
claim in $k_{\mathrm{piv}}$.
\end{restatedconjecture}

\emph{Status:} theorem at $T=1$ (Proposition~\ref{prop:bundle}(ii)); at the C2-falsifying
instance the refined scan is consistent with it (argmax $(0.8625,0)$); at risk from the
hunter's 6/8 genericity report if those violations survive $k_{\mathrm{type}}$-grid
refinement --- open. \emph{Falsifier:} any A0--A5 instance and $T$ where a fine
$k_{\mathrm{type}}$ scan certifies
$\sup_{k_{\mathrm{type}}}W(k_{\mathrm{type}},0)<W(k_{\mathrm{type}}',k_{\mathrm{piv}}')$ for
some $k_{\mathrm{piv}}'>0$. \emph{Obstruction (unchanged from the original C2):} for $T\ge 2$,
$k_{\mathrm{piv}}$ feeds back into $a_{\mathrm{mom}}$ through the continuation values
$\Pi,U$, so the clean factorization used in part (ii) fails state-by-state --- and O1 shows
the feedback can be strictly beneficial, so no monotone (Topkis-style) comparative static in
$k_{\mathrm{piv}}$ can hold pointwise; any proof of C2$'$ must compare optima, not paths.

\begin{remark}[R5 --- reading]\label{rem:reading}
Part (i)'s endpoints are the two channels' pure failure modes derived from primitives: A1
makes the \emph{pooled} commit option dominated (masking destroys the type news needed to
ignite), and A2$+$A3 make the \emph{root} mover provably inframarginal under full disclosure
(leakage destroys the perceived pivotality needed to move first). Part (ii) shows neither
failure is about \emph{balancing} the channels: given two instruments, the planner sets each
at its own best corner. The interior $k^*$ exists only because one physical k-anonymity dial
yokes the two --- exactly the v9 S11 claim, here as a theorem at $T=1$. Beyond one round the
two extensions now part ways: the bundled-interior extension C1 remains an open conjecture
(Conjecture~\ref{conj:interiorT}), while the decoupled-corner extension C2 is falsified
(Remark~\ref{rem:falsifiedC2}) --- at $T\ge 2$ positive leakage can strictly raise welfare
through the continuation channel (Observation~\ref{obs:leakagehelps}), so the corner result
is a one-round theorem plus an anchor-calibration numerical fact, not a general-horizon
truth.
\end{remark}
